% !PRINTED 595
\weberpage{595}
%<<CONT
und aus (5) ergibt sich durch Quadrieren mit Benutzung von (4):
\[
\pi^2 \frac{4\psi(\boldsymbol{s})^4 - \lambda\varphi(\boldsymbol{s})^2\psi(\boldsymbol{s})^2 + 4\psi(\boldsymbol{s})^4}{4 - \lambda \boldsymbol{s}^2 + 4\boldsymbol{s}^4} = [\varphi'(\boldsymbol{s})\psi(\boldsymbol{s}) - \varphi'(\boldsymbol{s})\psi(\boldsymbol{s})]^2. % [sic] both repeated numerator term and repeated right-hand factor are printed; (5) and the next display give the intended complementary terms
\]
Da rechts eine ganze Funktion von $\boldsymbol{s}$ steht, so muß auf der linken Seite der Zähler durch den Nenner teilbar sein (was sich auch aus der Betrachtung der Nullpunkte ergibt), und es folgt:
\[
\pi^2 \Phi = [\varphi'(\boldsymbol{s})\psi(\boldsymbol{s}) - \varphi(\boldsymbol{s})\psi'(\boldsymbol{s})]^2,
\]
wo $\Phi$ eine ganze Funktion von $\boldsymbol{s}$ ist, deren Koeffizienten ganze Zahlen in $\mathfrak{L}$ sind. Es muß also $\Phi$ das Quadrat einer ganzen Funktion sein, deren Koeffizienten gleichfalls ganze algebraische Zahlen sind (Bd. I, § 2; Bd. II, § 159), und es ergibt sich daraus für ein variables $x$:
\weq{6}{\varphi'(x)\psi(x) - \psi'(x)\varphi(x) \equiv 0 \ (\mathrm{mod}\ \pi).}
Hiernach läßt sich beweisen, daß die Koeffizienten $A_1, A_3, A_5, \ldots, A_{p-2}$ alle durch $\pi$ teilbar sind.

Von $A_1$ wissen wir das schon. Nehmen wir also an, es sei bewiesen für
\[
A_1, A_3, \ldots, A_{2\nu-1} \quad (2\nu+1 < p),
\]
so folgt:
\[
\begin{aligned}
\varphi(x) &\equiv A_{2\nu+1}x^{2\nu+1} + \cdots\\
\varphi'(x) &\equiv (2\nu+1)A_{2\nu+1}x^{2\nu} + \cdots \qquad (\mathrm{mod}\ \pi)\\
\psi(x) &\equiv 1 + A_{p-2}x^2 + \cdots\\
\psi'(x) &\equiv 2A_{p-2}x + \cdots
\end{aligned}
\]
Suchen wir also in (6) den Koeffizienten von $x^{2\nu}$, so folgt:
\[
(2\nu+1)A_{2\nu+1} \equiv 0 \ (\mathrm{mod}\ \pi).
\]
Da $2\nu+1 < p$ und $p$ eine Primzahl ist, so ist $2\nu+1$ relativ prim zu $\pi$, und es folgt:
\weq{7}{A_{2\nu+1} \equiv 0 \ (\mathrm{mod}\ \pi).}
Setzen wir in (4) für $u$ einen der Werte $2\alpha\xi/\mu$, § 158, (9), so wird $\boldsymbol{s}(u)$ eine Wurzel der Ideal-Teilungsgleichung § 158, (10), und es folgt:
\[
\pm \boldsymbol{s}(\pi u) \equiv [\boldsymbol{s}(u)]^p \ (\mathrm{mod}\ \pi),
\]
und wenn wir ein Quadrat erheben:
\weq{8}{\boldsymbol{s}(\pi u)^2 \equiv [\boldsymbol{s}(u)]^{2p} \ (\mathrm{mod}\ \pi).}
Insbesondere ist also
\weq{9}{\boldsymbol{s}(u)^2 \equiv \boldsymbol{s}(u)^{2p} \ (\mathrm{mod}\ \pi),}
wenn
\weq{10}{\pi \equiv \pm 1 \ (\mathrm{mod}\ \mathfrak{m})}
%CONT>>
